\begin{abstract}
“板凳龙”由多节板凳通过把手首尾相连，运动过程中既受到螺线轨迹约束，又受到相邻把手距离、板凳实体宽度和调头空间边界的共同限制。本文围绕盘入、调头和盘出过程，建立以\textbf{等距螺线方程}、\textbf{极坐标弧长积分}、\textbf{逐节几何递推}、\textbf{有限宽度碰撞判定}和\textbf{速度比例缩放}为核心的统一模型，分别求解各时刻位置速度、无碰撞终止时刻、最小螺距、调头路径以及最大安全头速。

\textbf{针对问题一，}首先用\textbf{等距螺线方程}刻画半径随极角线性变化的盘入轨迹，再由极坐标参数曲线的弧长微元积分得到\textbf{螺线弧长函数}，据此反解龙头前把手的角参数；随后利用\textbf{相邻把手距离约束}逐节递推全队位置，并用\textbf{差分法}计算速度。模型生成了0--300 s内全部224个把手的位置和速度，位置表规模为 \textbf{449×302}，速度表规模为 \textbf{225×302}，相邻把手距离最大误差小于 \textbf{1.0e-11 m}。

\textbf{针对问题二，}将每节板凳视为有限宽度矩形，依据把手连线确定矩形轴线和法向，并采用\textbf{分离轴定理}判断非相邻板凳是否相交。再定义碰撞判定函数，通过\textbf{扫描}确定临界区间，并用\textbf{二分法}逼近首次碰撞时刻。无碰撞终止时刻为 \textbf{412.474 s}，首次碰撞板对为 \textbf{[0,8]}，终止状态下全队速度约在 \textbf{0.973--1.000 m/s} 之间。

\textbf{针对问题三，}在问题二碰撞判定思想的基础上，以龙头附近外侧角点到前方板凳轴线的最小距离作为接触判据。采用\textbf{变步长搜索}：以0.55 m为初始螺距逐步缩小，并在出现接触风险后减小步长继续逼近。满足盘入调头空间且不发生接触风险的最小螺距为 \textbf{0.45 m}。

\textbf{针对问题四，}在螺距1.7 m、调头空间半径4.5 m的条件下，以盘入螺线与调头圆的交点作为调头入点，并令盘出螺线与盘入螺线关于中心对称。根据\textbf{两圆弧相切模型}约束两段圆弧分别与盘入、盘出螺线相切、两圆弧相切以及半径比为2:1的几何关系，从而确定调头路径。两段圆弧半径分别为 \textbf{3.005 m} 和 \textbf{1.503 m}，调头路径长度为 \textbf{13.621 m}；同时输出调头前后100 s内每秒全队位置和速度。

\textbf{针对问题五，}在问题四固定路径上计算单位龙头速度下的全队速度矩阵。由于路径固定时各把手速度随龙头恒速成比例缩放，采用\textbf{单位头速归一化}得到最大速度放大倍数，并反推出龙头速度上界。单位头速下最大速度放大倍数为 \textbf{1.422}，限制把手为 \textbf{第1节龙身}，对应龙头最大恒定速度为 \textbf{1.406 m/s}。

\keywords{等距螺线\quad 几何递推\quad 碰撞判定\quad 二分搜索}
\end{abstract}
